\chapter[SQL Avanzado]{\emj{🚀}~SQL Avanzado}

\section[CTEs (Common Table Expressions)]{\emj{🧠}~CTEs (Common Table Expressions)}

Una CTE es una subconsulta con nombre definida antes del \texttt{SELECT} principal usando la cláusula \texttt{WITH}.
El objetivo principal es legibilidad: en lugar de anidar subconsultas dentro de otras subconsultas (que se vuelven ilegibles rápidamente), se definen como bloques nombrados que se leen de arriba hacia abajo como un programa.

Una CTE puede referenciarse múltiples veces en la misma query, lo que evita duplicación.
En PostgreSQL, una CTE por defecto se \textbf{materializa}: se ejecuta una vez y el resultado se almacena en memoria.
Esto puede ser una ventaja (si se referencia múltiples veces) o una desventaja (si el planificador no puede optimizarla junto con el resto de la query).
Con \texttt{NOT MATERIALIZED}, le indicas al planificador que trate la CTE como una vista expandida en línea.

Las \textbf{CTEs recursivas} usan \texttt{WITH RECURSIVE} y se auto-referencian.
Son la herramienta estándar en SQL para recorrer estructuras jerárquicas (árboles organizacionales, categorías anidadas, grafos).
La recursión termina cuando la parte recursiva no retorna filas.

\begin{itemize}
  \item CTE regular: alias de subconsulta reutilizable, mejora legibilidad
  \item CTE recursiva (\texttt{WITH RECURSIVE}): auto-referencia para jerarquías y grafos
  \item \texttt{MATERIALIZED}: fuerza ejecución y caché del resultado
  \item \texttt{NOT MATERIALIZED}: permite al planificador optimizar junto con la query externa
\end{itemize}

\section[Window Functions]{\emj{📊}~Window Functions}

Las window functions son una de las características más poderosas de SQL moderno.
Permiten calcular valores en relación a un conjunto de filas (la ``ventana'') \textbf{sin colapsar el resultado} como hace \texttt{GROUP BY}.
Cada fila del resultado mantiene su identidad y adicionalmente puede ver valores de sus filas vecinas.

El caso clásico que las necesita: calcular el salario de cada empleado y también el promedio del salario de su departamento en la misma fila.
Con \texttt{GROUP BY} esto requeriría un JOIN adicional; con window functions es una sola expresión.

\begin{teoriabox}{Anatomía de una Window Function}
\texttt{función() OVER (PARTITION BY ... ORDER BY ... ROWS/RANGE BETWEEN ...)}
\begin{itemize}
  \item \textbf{función():} \texttt{ROW\_NUMBER, RANK, DENSE\_RANK, LAG, LEAD, FIRST\_VALUE, LAST\_VALUE, SUM, AVG, COUNT}
  \item \textbf{PARTITION BY:} divide las filas en grupos independientes (como GROUP BY pero sin colapsar filas). Cada partición tiene su propia ventana.
  \item \textbf{ORDER BY:} define el orden dentro de cada partición. Obligatorio para funciones que dependen de posición relativa (LAG, LEAD, running totals).
  \item \textbf{ROWS BETWEEN:} delimita exactamente cuántas filas forman la ventana para cada fila actual.
    \texttt{UNBOUNDED PRECEDING}: desde el inicio de la partición.
    \texttt{CURRENT ROW}: la fila actual.
    \texttt{UNBOUNDED FOLLOWING}: hasta el final de la partición.
\end{itemize}
\textbf{ROW\_NUMBER} asigna un número único siempre; \textbf{RANK} salta números con empates (1,1,3); \textbf{DENSE\_RANK} no salta (1,1,2).
\textbf{LAG(col, n)} devuelve el valor de \texttt{col} de n filas \emph{antes}; \textbf{LEAD(col, n)} de n filas \emph{después}.
\end{teoriabox}

\section[EXPLAIN ANALYZE]{\emj{🔍}~EXPLAIN ANALYZE}

\texttt{EXPLAIN} y \texttt{EXPLAIN ANALYZE} son las herramientas más importantes para entender por qué una query es lenta y cómo optimizarla.
\texttt{EXPLAIN} muestra el plan sin ejecutar la query (estimaciones).
\texttt{EXPLAIN ANALYZE} ejecuta la query y muestra tiempos reales junto con las estimaciones, permitiendo comparar.

El plan se muestra como un árbol de nodos.
Cada nodo representa una operación; el costo se acumula de hojas a raíz.
El nodo más interno (más indentado) se ejecuta primero.

\begin{itemize}
  \item \textbf{Seq Scan:} lee toda la tabla secuencialmente. Normal para tablas pequeñas o cuando se retorna más del ~5-10\% de las filas. Preocupante en tablas grandes con filtros selectivos.
  \item \textbf{Index Scan:} usa un índice para localizar filas específicas. Eficiente para alta selectividad (pocos resultados de muchas filas).
  \item \textbf{Bitmap Index Scan + Bitmap Heap Scan:} acumula punteros del índice, los ordena y accede al heap en orden. Más eficiente que múltiples Index Scans.
  \item \textbf{Hash Join:} construye una hash table del lado más pequeño y sondea el grande. Eficiente para joins de tablas medianas/grandes.
  \item \textbf{Nested Loop:} itera el lado externo y busca en el interno. Eficiente cuando el lado interno es pequeño o indexado.
  \item \textbf{Merge Join:} ambas tablas ordenadas; avanza en paralelo. Eficiente cuando los datos ya están ordenados.
  \item \textbf{cost=inicio..total:} unidades de costo abstractas (no segundos). El inicio es el costo antes de la primera fila; el total es el costo de retornar todas las filas.
  \item \textbf{rows:} estimación de filas. Si difiere mucho del \texttt{actual rows}, las estadísticas están desactualizadas. Ejecuta \texttt{ANALYZE tabla} para actualizarlas.
\end{itemize}

\begin{lstlisting}[caption={CTEs --- regular y recursiva}]
-- CTE regular: ventas por región con total acumulado
WITH ventas_region AS (
    SELECT region, SUM(monto) AS total
    FROM   ventas
    WHERE  fecha >= '2024-01-01'
    GROUP BY region
)
SELECT region, total,
       ROUND(100.0 * total / SUM(total) OVER (), 2) AS pct
FROM   ventas_region
ORDER BY total DESC;

-- CTE recursiva: jerarquía de empleados (árbol)
WITH RECURSIVE jerarquia AS (
    SELECT id, nombre, id_jefe, 0 AS nivel
    FROM   empleados WHERE id_jefe IS NULL

    UNION ALL

    SELECT e.id, e.nombre, e.id_jefe, j.nivel + 1
    FROM   empleados e
    JOIN   jerarquia j ON e.id_jefe = j.id
)
SELECT nivel, nombre FROM jerarquia ORDER BY nivel, nombre;
\end{lstlisting}

\begin{lstlisting}[caption={Window Functions}]
-- ROW_NUMBER, RANK, DENSE_RANK
SELECT nombre, depto, salario,
       ROW_NUMBER() OVER (PARTITION BY depto ORDER BY salario DESC) AS fila,
       RANK()       OVER (PARTITION BY depto ORDER BY salario DESC) AS rango,
       DENSE_RANK() OVER (PARTITION BY depto ORDER BY salario DESC) AS rango_denso
FROM empleados;

-- LAG y LEAD: comparar con fila anterior/siguiente
SELECT fecha, ventas,
       LAG(ventas)  OVER (ORDER BY fecha) AS ventas_dia_anterior,
       ventas - LAG(ventas) OVER (ORDER BY fecha) AS delta
FROM reporte_diario;

-- Running total
SELECT fecha, monto,
       SUM(monto) OVER (ORDER BY fecha ROWS UNBOUNDED PRECEDING) AS acumulado
FROM ventas;
\end{lstlisting}

\begin{lstlisting}[caption={EXPLAIN ANALYZE}]
EXPLAIN ANALYZE
SELECT e.nombre, d.nombre
FROM   empleados e
JOIN   departamentos d ON e.id_depto = d.id
WHERE  e.salario > 50000;

-- Salida tipica:
-- Hash Join  (cost=12.50..34.20 rows=50 width=200)
--   -> Seq Scan on empleados
--        Filter: (salario > 50000)
--   -> Hash
--        -> Seq Scan on departamentos
-- Si ves "Seq Scan" en tabla grande = considera indice en esa columna
\end{lstlisting}

\begin{entrevistabox}
\begin{enumerate}
  \item ¿Cuál es la diferencia entre \texttt{ROW\_NUMBER()}, \texttt{RANK()} y \texttt{DENSE\_RANK()}?
  \item ¿Cuándo usarías una CTE recursiva? Da un ejemplo de caso de uso.
  \item ¿Cómo interpretas un nodo \texttt{Seq Scan} en \texttt{EXPLAIN ANALYZE}?
\end{enumerate}
\end{entrevistabox}
